\chapter{Floaters and Vitreous Detachment}
\index{Floaters and Vitreous Detachment}
\label{sec:Floaters and Vitreous Detachment}

\section{Overview}
\begin{itemize}[noitemsep]
  \item Floaters are small shadows, specks, threads, or cobweb-like shapes that drift across the field of vision as the eye moves
  \item Posterior vitreous detachment (PVD) is the age-related separation of the vitreous gel from the retina, and is the most common cause of sudden onset of floaters
  \item PVD occurs in most people by age 80 and is usually a normal part of aging rather than a disease
\end{itemize}
\section{Causes}
Floaters and vitreous detachment happen because of age-related liquefaction (syneresis) and shrinkage of the vitreous gel. As the gel shrinks, it pulls away from the retina; the strands and debris cast shadows on the retina that are perceived as floaters. Secondary causes include trauma, high myopia, intraocular inflammation (uveitis), bleeding into the vitreous (vitreous hemorrhage), and previous eye surgery such as cataract extraction.
\section{Risk Factors}
\begin{itemize}[noitemsep]
  \item Age over 50 years
  \item Nearsightedness (high myopia)
  \item Eye trauma or prior eye surgery
  \item Intraocular inflammation (uveitis)
  \item Diabetes with retinal disease
  \item Family history of early vitreous detachment
\end{itemize}
\section{Signs and Symptoms}
\subsection{Common symptoms}
\begin{itemize}[noitemsep]
  \item Seeing spots, specks, threads, or cobweb shapes that drift when the eye moves
  \item Flashes of light (photopsia), often in the peripheral vision, especially in dim lighting
  \item Sudden shower of many new floaters
  \item A shadow or "curtain" across part of the vision (a retinal emergency)
\end{itemize}
\subsection{Emergency warning signs}
\begin{itemize}[noitemsep]
  \item Sudden appearance of many new floaters
  \item Flashes of light that persist
  \item A dark curtain or veil covering part of the visual field
  \item Sudden loss of vision --- these may indicate a retinal tear or detachment and require immediate eye examination
\end{itemize}
\section{Progression and Stages}
The vitreous detachment is a single event that usually occurs over weeks to months. Floaters are most noticeable when they first appear and typically become less bothersome over time as the debris settles and the brain adapts. The critical danger window is in the weeks after PVD begins, when vitreous traction can cause a retinal tear.
\begin{itemize}[noitemsep]
  \item Vitreous syneresis (gel liquefaction) precedes detachment
  \item Complete PVD separates the gel from the retina
  \item Incomplete or anomalous PVD may leave traction on the retina, leading to tears, hemorrhage, or detachment
\end{itemize}
\section{Complications}
\begin{itemize}[noitemsep]
  \item Retinal tear (tractional)
  \item Retinal detachment
  \item Vitreous hemorrhage
  \item Epiretinal membrane formation
  \item Persistent visually disturbing floaters affecting daily function
\end{itemize}
\section{Diagnosis}
\begin{itemize}[noitemsep]
  \item Dilated fundus examination with scleral depression to detect tears
  \item Slit-lamp biomicroscopy to inspect the vitreous
  \item B-scan ultrasound when the vitreous is too hazy to see the retina
  \item Optical coherence tomography (OCT) if macular pathology is suspected
  \item Urgent referral when new floaters, flashes, or visual field loss are present
\end{itemize}
\section{Prevention}
\begin{itemize}[noitemsep]
  \item Posterior vitreous detachment cannot be prevented; it is a normal aging change
  \item Prompt examination for sudden onset of floaters or flashes to allow early treatment of any tear
  \item Protect the eyes from trauma
  \item Control underlying conditions such as diabetes
\end{itemize}
\section{Treatment Options}
\begin{itemize}[noitemsep]
  \item No treatment is needed for uncomplicated PVD and benign floaters; symptoms usually improve with time
  \item Laser retinopexy or cryotherapy to seal retinal tears
  \item Vitrectomy surgery for dense, persistent, or visually disabling floaters
  \item Surgery (scleral buckle, vitrectomy) if retinal detachment develops
  \item Regular follow-up monitoring during the acute phase
\end{itemize}
\section{Living with It}
\begin{itemize}[noitemsep]
  \item Most floaters become less noticeable as the brain learns to ignore them
  \item Follow up as directed, especially in the weeks after symptoms begin
  \item Seek urgent care if new flashes, more floaters, or a shadow appear
  \item Stay informed about your condition
\end{itemize}
\section{Outlook}
The outlook is excellent for uncomplicated posterior vitreous detachment; most people experience harmless floaters that fade in prominence. Prompt treatment of retinal tears prevents most detachments. When retinal detachment occurs, timely surgery usually restores vision, though outcomes depend on how quickly treatment is delivered.
